% Source: Wald GR, physical PDF p. 18 (printed p. 5).
% Coverage: Figure 1.1, causal structure of spacetime in prerelativity physics.
% Figures/tables/footnotes: Figure only; no tables or footnotes.
% Status: source-reviewed and compile-clean.
% Uncertainties: none.

\begingroup
\renewcommand{\thefigure}{1.1}
\begin{figure}[htbp]
  \centering
  \begin{tikzpicture}[>=Stealth, line cap=round, line join=round]
    \draw[->, thick] (0,4.40) -- (0,5.05) node[above] {time};
    \draw[->, thick] (0,4.40) -- (0.85,4.40) node[right] {space};
    \draw[->, thick] (0,4.40) -- (-0.34,4.06);

    \draw[thick] (0.65,2.20) -- (5.65,2.20) -- (6.90,3.75) --
      (1.90,3.75) -- cycle;
    \fill (3.40,2.94) circle (1.7pt) node[left=4pt] {\(p\)};
    \node at (4.15,4.35) {Future};
    \node at (3.10,1.60) {Past};
    \node[anchor=west] at (6.25,4.10) {Simultaneous};
    \draw[->, thick, bend right=18] (6.36,3.97) to (5.88,3.10);
  \end{tikzpicture}
  \caption{A diagram showing the causal structure of spacetime in prerelativity
  physics. Given an event \(p\), all other events in spacetime either are to the future
  of \(p\), to the past of \(p\), or simultaneous with \(p\). The simultaneous events
  form a three-dimensional surface in spacetime.}
  \label{fig:1.1}
\end{figure}
\endgroup
